%------------------------------------------------------------------------------%
\begin{question}
  Diagonalize ortogonalmente a matriz
  \[
    A =
    \begin{pmatrix}
      2 & 1 & 0 \\
      1 & 2 & 0 \\
      0 & 0 & 4
    \end{pmatrix}
  \]
  Determine as matrizes $P$ e $D$ tais que
  \(
    A = P D P^t
  \)
\end{question}

%------------------------------------------------------------------------------%
\begin{resolution}{Autovalores}
  Polinômio característico
  \begin{align*}
    \onslide<+->{p(\lambda)}
    \onslide<+->{= \det(A-\lambda I) }
    \onslide<+->{& = \begin{vmatrix}
      2-\lambda & 1         & 0         \\
      1         & 2-\lambda & 0         \\
      0         & 0         & 4-\lambda
    \end{vmatrix}}
    \\
    \onslide<+->{& = (-1)^{3+3} (4-\lambda)
      \begin{vmatrix}
        2-\lambda & 1         \\
        1         & 2-\lambda
      \end{vmatrix}}
    \\
    \onslide<+->{& = (4-\lambda) \left[ (2-\lambda)^2 - 1            \right]} \\
    \onslide<+->{& = (4-\lambda) \left[ 4 - 4\lambda + \lambda^2 - 1 \right]} \\
    \onslide<+->{& = (4-\lambda) \left( \lambda^2 - 4\lambda + 3     \right)}
  \end{align*}
\end{resolution}

%------------------------------------------------------------------------------%
\begin{resolution}{Autovalores}
  \[
    \onslide<+->{(4-\lambda) \left( \lambda^2 - 4\lambda + 3 \right) = 0}
  \]

  \[
    \onslide<+->{\lambda_1 = 4}
    \onslide<+->{\qquad\text{ ou }\qquad
    \lambda^2 - 4\lambda + 3 = 0}
  \]

  \[
    \onslide<+->{\Delta}
    \onslide<+->{= b^2 - 4ac}
    \onslide<+->{= (-4)^2 - 4 \times 1 \times 3}
    \onslide<+->{= 4}
    \qquad\qquad
    \onslide<+->{\sqrt{\Delta}}
    \onslide<+->{= \sqrt{4}}
    \onslide<+->{= 2}
  \]

  \[
    \onslide<+->{\lambda_2}
    \onslide<+->{= \frac{-b+\sqrt{\Delta}}{2a}}
    \onslide<+->{= \frac{4+2}{2}}
    \onslide<+->{= 3}
  \]

  \[
    \onslide<+->{\lambda_3}
    \onslide<+->{= \frac{-b-\sqrt{\Delta}}{2a}}
    \onslide<+->{= \frac{4-2}{2}}
    \onslide<+->{= 1}
  \]
\end{resolution}

%------------------------------------------------------------------------------%
\begin{resolution}{Autovetores para $\lambda_1=4$}
\begin{columns}
\column{0.4\textwidth}
  \begin{align*}
    \onslide<+->{(A-4I)\vec{x}&=\vec{0}}
    \\
    \onslide<+->{\begin{bmatrix}
      -2 & 1  & 0 \\
      1  & -2 & 0 \\
      0  & 0  & 0
    \end{bmatrix}
    \vetortri{x}{y}{z} & = \vetortri{0}{0}{0}}
  \end{align*}
  \vspace{-\baselineskip}
  \begin{align*}
    \onslide<+->{-2x +  y & = 0 \\}
    \onslide<+->{y & = 2x}
    \\[2mm]
    \onslide<+->{x - 2y & =}
    \onslide<+->{x - 4x =}
    \onslide<+->{-x = 0} \\
    \onslide<+->{y & = 0}
  \end{align*}
\column{0.4\textwidth}
  \onslide<+->{Escolhendo $z=1$}
  \onslide<+->{\[
    \vec{v}_1 = \vetortri{0}{0}{1}
  \]}
  \onslide<+->{\[
    E(4) = \ger \left\{ \vetortri{0}{0}{1} \right\}
  \]}
\end{columns}
\end{resolution}

%------------------------------------------------------------------------------%
\begin{resolution}{Normalização de $\vec{v}_1$}
  Norma de
  \(
    \vec{v}_1 = \vetortri{0}{0}{1}
  \)
  \pause
  \[
    \norm{\vec{v}_{2}}
    \pause
    = \sqrt{0^2 +0^2 + 1^2}
    \pause
    = 1
  \]

  \pause
  Autovetor unitário
  \[
    \pause
    \vec{u}_1
    \pause
    = \vetortri{0}{0}{1}
  \]
\end{resolution}

%------------------------------------------------------------------------------%
\begin{resolution}{Autovetores para $\lambda_2=3$}
\begin{columns}
\column{0.4\textwidth}
  \begin{align*}
    \onslide<+->{(A-3I)\vec{x}&=\vec{0}}
    \\
    \onslide<+->{\begin{bmatrix}
      -1 & 1  & 0 \\
      1  & -1 & 0 \\
      0  & 0  & 1
    \end{bmatrix}
    \vetortri{x}{y}{z} & = \vetortri{0}{0}{0}}
  \end{align*}
  \onslide<+->{\begin{align*}
    -x + y & = 0 \\
    z      & = 0
  \end{align*}}
\column{0.4\textwidth}
  \onslide<+->{Escolhendo $x=1$}
  \onslide<+->{\[
    y = 1
  \]}
  \onslide<+->{\[
    \vec{v}_2 = \vetortri{1}{1}{0}
  \]}
  \onslide<+->{\[
    E(3) = \ger \left\{ \vetortri{1}{1}{0} \right\}
  \]}
\end{columns}
\end{resolution}

%------------------------------------------------------------------------------%
\begin{resolution}{Normalização de $\vec{v}_2$}
  Norma de
  \(
    \vec{v}_2 = \vetortri{1}{1}{0}
  \)
  \pause
  \[
    \norm{\vec{v}_{2}}
    \pause
    = \sqrt{1^2 +1^2 + 0^2}
    \pause
    = \sqrt{2}
  \]

  \pause
  Autovetor unitário
  \[
    \pause
    \vec{u}_2
    \pause
    = \frac{1}{\sqrt{2}}
      \vetortri{1}{1}{0}
  \]
\end{resolution}

%------------------------------------------------------------------------------%
\begin{resolution}{Autovetores para $\lambda_3=1$}
\begin{columns}
\column{0.4\textwidth}
  \begin{align*}
    \onslide<+->{(A-I)\vec{x}&=\vec{0}}
    \\
    \onslide<+->{\begin{bmatrix}
      1 & 1 & 0 \\
      1 & 1 & 0 \\
      0 & 0 & 3
    \end{bmatrix}
    \vetortri{x}{y}{z} & = \vetortri{0}{0}{0}}
  \end{align*}
  \onslide<+->{\begin{align*}
    x+y & = 0 \\
    z   & = 0
  \end{align*}}
\column{0.4\textwidth}
  \onslide<+->{Escolhendo $x=1$}
  \onslide<+->{\[
    y = -1
  \]}
  \onslide<+->{\[
    \vec{v}_3 = \vetortri{1}{-1}{0}
  \]}
  \onslide<+->{\[
    E(1) = \ger \left\{ \vetortri{1}{-1}{0} \right\}
  \]}
\end{columns}
\end{resolution}

%------------------------------------------------------------------------------%
\begin{resolution}{Normalização de $\vec{v}_3$}
  Norma de
  \(
    \vec{v}_3 = \vetortri{1}{-1}{0}
  \)
  \pause
  \[
    \norm{\vec{v}_{3}}
    \pause
    = \sqrt{1^2 +(-1)^2 + 0^2}
    \pause
    = \sqrt{2}
  \]

  \pause
  Autovetor unitário
  \[
    \pause
    \vec{u}_3
    \pause
    = \frac{1}{\sqrt{2}}
      \vetortri{1}{-1}{0}
  \]
\end{resolution}

%------------------------------------------------------------------------------%
\begin{resolution}{Base Ortonormal}
  Autovetores unitários
  \[
    \vec{u}_{1} = \vetortri{0}{0}{1}                     \qquad
    \vec{u}_{2} = \frac{1}{\sqrt{2}} \vetortri{1}{1}{0}  \qquad
    \vec{u}_{3} = \frac{1}{\sqrt{2}} \vetortri{1}{-1}{0}
  \]
\end{resolution}

%------------------------------------------------------------------------------%
\begin{resolution}{Diagonalização}
  Organizamos os autovetores nas colunas de $P$
  \[
    P =
    \begin{pmatrix}
      \dfrac{1}{\sqrt{2}}  & \dfrac{1}{\sqrt{2}} & 0 \\[0.3cm]
      -\dfrac{1}{\sqrt{2}} & \dfrac{1}{\sqrt{2}} & 0 \\[0.3cm]
      0                    & 0                   & 1
    \end{pmatrix}
  \]

  Organizamos os autovalores na diagonal de $D$
  \[
    D =
    \begin{pmatrix}
      1 & 0 & 0 \\
      0 & 3 & 0 \\
      0 & 0 & 4
    \end{pmatrix}
  \]
\end{resolution}

%------------------------------------------------------------------------------%
